\begin{textsample}{NEG Dim 1 – 2016\_c5 – Score 9.00 – 2016\_c5/enade\_2016\_zoot\_f32beec87b.txt}  \label{ex:f1_neg_366}
CAPÍTULO 5 – DISTRIBUIÇÃO DOS CONCEITOS 5.1 - PANORAMA NACIONAL DA DISTRIBUIÇÃO DOS CONCEITOS O Conceito Enade é um indicador de qualidade que avalia o desempenho dos estudantes a partir dos resultados do Enade.
Os valores possíveis vão de 1 ( pior situação ) a 5 ( melhor situação ).
As Unidades de Observação com apenas um ou sem nenhum concluinte participante não obtêm o Conceito Enade, ficando Sem Conceito ( SC ).
A Tabela 5.1 apresenta a quantidade e distribuição de cursos de Zootecnia participantes do Enade/2016, por faixa de conceito e Grande Região.
Enfatiza -se, mais uma vez, que a diferença entre os cursos tabulados neste capítulo e no capítulo 2 corresponde aos cursos sem conceito, ou seja àqueles, em princípio, sem alunos concluintes que houvessem participado da prova ou àqueles que tivessem tão somente um aluno concluinte19.
Observando -se os dados da Tabela 5.1, nota -se que, dos 93 cursos participantes, 38 ( 40,9 \% ) classificaram -se com conceito 3, o valor modal.
Esse foi também o conceito modal em três regiões : Nordeste ( 38,5 \% ), Sudeste ( 59,1 \% ) e Sul ( 50,0 \% ).
Nas demais regiões, a moda foi o conceito 2, atribuído para 54,5 \% dos cursos da região Norte e para 50,0 \% dos cursos da região Centro-Oeste.
O conceito 2 foi o segundo mais frequente em nível nacional ( 29,0 \%, correspondendo a 27 cursos ) e o conceito 4, o terceiro ( 19,4 \%, correspondendo a 18 cursos ).
Houve, ainda, quatro cursos ( 4,3 \% ) que receberam conceito 5 e cinco cursos ( 5,4 \% ) que receberam conceito 1.
Dos 93 cursos de Zootecnia, somente um ficou sem conceito ( SC ). 19 Estes cursos com somente um concluinte não foram considerados no capítulo 2.
A região Norte participou com 11 cursos ou 11,8 \% do total nacional.
Desses, seis cursos receberam o conceito 2, o que equivale a 54,5 \% do total regional, além de ser o conceito modal na região, como já comentado.
A dois cursos atribuíram -se conceitos 1 e 5 ( um curso, \textbf{cada} ) e a três cursos ( 27,3 \% ), conceito 3.
Nenhum curso ficou sem conceito ( SC ) ou recebeu conceito 4 nessa região.
A região Nordeste participou com 26 cursos ou 28,0 \% do total nacional.
Nessa região, dez cursos ( 38,5 \% em termos regionais ) obtiveram conceito 3, o conceito modal para a região, como já comentado.
Os conceitos 1 e 2 foram atribuídos, respectivamente, a três e a sete cursos ( 11,5 \% e 26,9 \%, respectivamente ).
Ainda nessa região, cinco cursos ( 19,2 \% ) foram avaliados com o conceito 4.
Apenas um curso ficou sem conceito ( SC ), enquanto nenhum recebeu conceito 5.
Dos 22 cursos participantes da região Sudeste ( 23,7 \% do total nacional ), o conceito 3 foi o conceito modal, como já anteriormente assinalado, tendo sido obtido por 13 cursos ( 59,1 \% ).
O conceito 2 foi atribuído a três cursos ( 13,6 \% ) e o conceito 4 a cinco cursos ( 22,7 \% ).
Receberam o conceito 5, um único curso ( 4,5 \% ).
Nessa região, nenhum curso ficou sem conceito ( SC ) ou com conceito 1.
Os 16 cursos da região Sul corresponderam a 17,2 \% do total nacional.
Como já indicado, houve predominância do conceito 3, atribuído a 50,0 \% dos cursos da região ( conceito modal ), ou seja, atribuído a oito dos 16 cursos participantes na região Sul.
O conceito 2 foi atribuído a dois cursos ( 12,5 \% ), o conceito 4, a cinco cursos ( 31,3 \% ) e o conceito 5, a um curso ( 6,3 \% ).
Nenhum curso recebeu conceito 1 ou ficou sem conceito ( SC ) na região Sul.
Já dos 18 cursos participantes na região Centro-Oeste ( 19,4 \% do total nacional ), nove receberam conceito 2, o conceito modal ( 50,0 \% ), como já fora destacado.
A dois cursos foram atribuídos conceitos 1 e 5 ( um curso, \textbf{cada} ).
Os conceitos restantes, 3 e 4, foram atribuídos, respectivamente, a quatro ( 22,2 \% ) e três ( 16,7 \% ) cursos.
Nenhum curso ficou sem conceito ( SC ).
O Gráfico 5.1 apresenta a distribuição do Conceito Enade segundo Grande Região : poligonais mais à esquerda representam Grandes Regiões com pior distribuição, e poligonais mais à direita, Grandes Regiões com melhores conceitos.
A região Norte ( linha azul ) apresenta os cursos com conceitos mais baixos, ao passo que as regiões Sudeste ( linha vermelha ) e Sul ( linha roxa ) apresentam as distribuições com valores maiores.
Cursos em IES situadas nas regiões Nordeste ( linha verde ) e Centro-Oeste ( linha preta ) apresentam uma situação intermediária. 5.2 - CONCEITOS POR CATEGORIA ADMINISTRATIVA E GRANDE REGIÃO A Tabela 5.2 apresenta a distribuição dos cursos participantes do Enade/2016 de Zootecnia, por Categoria Administrativa, de acordo com os conceitos por eles alcançados, segundo as Grandes Regiões brasileiras.
Dos 93 cursos participantes, 79 ( 84,9 \% ) eram ministrados em Instituições Públicas e 14 ( 15,1 \% ), em Instituições Privadas.
De acordo com as informações da Tabela 5.2, em termos nacionais, entre Instituições Públicas e Privadas, observa -se que, dos quatro cursos avaliados com conceito 5, três eram oferecidos em IES Públicas e um, em IES Privadas.
Dos 79 cursos participantes de IES Públicas, o conceito 3 foi o valor modal, atribuído a 36 cursos ( 45,6 \% ).
Entre os demais cursos de Instituições Públicas participantes, quatro cursos ( 5,1 \% ) foram avaliados com conceito 1, 18 cursos ( 22,8 \% da categoria ) obtiveram conceito 2, 18 cursos ( 22,8 \% da categoria ) obtiveram conceito 4 e, como já comentado, três cursos ( 3,8 \% ) receberam conceito 5.
Nesta categoria, nenhum curso ficou sem conceito ( SC ).
Na rede privada, o conceito modal foi 2, atribuído a nove cursos, o correspondente a 64,3 \% dos 14 cursos da categoria.
Entre os demais cursos participantes, um ( 7,1 \% ) recebeu conceito 1, e dois ( 14,3 \% ), conceito 3.
Os conceitos 4 e 5 foram atribuídos a um curso, \textbf{cada}.
Nesta Categoria Administrativa, nenhum curso ficou sem conceito ( SC ).
Na análise por região, observa -se que, na região Norte, as Instituições Privadas participaram com quatro cursos ( 36,4 \% do total regional ), a proporção mais elevada do que a registrada nas demais regiões brasileiras, dos quais um ficou com conceito 1 e três com conceito 2, conceito modal.
As Instituições Públicas participaram com sete cursos na região Norte ( 63,6 \% do total regional ).
Nota -se que tais cursos apresentaram distribuição bimodal, uma vez que as maiores frequências se concentraram em dois conceitos, os conceitos 2 e 3, atribuídos, \textbf{cada} qual, a três cursos ( 42,9 \% \textbf{cada} ).
Além disso, um curso alcançou o conceito máximo, 5.
Nenhum curso ficou sem conceito ( SC ) ou com conceito 1 ou 4 nessa combinação de categoria e região.
Na região Nordeste, a rede privada participou com somente um dos 26 cursos participantes, o equivalente a 3,8 \% do total da região.
Este curso recebeu conceito 2.
As Instituições Públicas da região Nordeste concentraram 25 cursos ( 96,2 \% do total da região ).
Desses, dez obtiveram conceito 3, o valor modal, seguidos por outros seis cursos, que receberam conceito 2.
Os conceitos 1 e 4 foram atribuídos, respectivamente, a três e a cinco cursos.
Apenas um curso ficou sem conceito ( SC ) nessa região, enquanto nenhum recebeu conceito 5.
Na região Sudeste, a proporção de cursos da rede privada foi de 22,7 \%, correspondendo a cinco dos 22 cursos participantes.
Nessa categoria e região, o conceito modal foi 3, atribuído a dois cursos.
Os três cursos restantes foram alocados aos conceitos 2, 4 e 5 ( um curso, \textbf{cada} ).
Nenhum curso ficou sem conceito ( SC ) ou com conceito 1.
Entre os 17 cursos oferecidos por Instituições Públicas na região Sudeste, o conceito modal foi 3, atribuído a 11 cursos.
Os demais cursos receberam os conceitos 2 ( dois cursos ) e 4 ( quatro cursos ).
Nenhum curso ficou sem conceito ( SC ) ou recebeu conceito 1 ou 5.
As Instituições Públicas concentraram todos os 16 cursos participantes da região Sul, 100,0 \% do total regional, a maior das proporções desta categoria nas regiões.
Desses, oito obtiveram conceito 3, o conceito modal.
Nesta combinação de Categoria Administrativa e Grande Região, os conceitos 2 e 4 foram atribuídos a dois e a cinco cursos, respectivamente.
Um curso recebeu conceito 5, e nenhum recebeu conceito 1 ou ficou sem conceito ( SC ).
As Instituições Públicas na região Sul não tiveram participação.
Na região Centro-Oeste, quatro dos 18 cursos participantes eram de Instituições Privadas ( 22,2 \% em termos regionais ), todos eles receberam conceito 2.
Quanto aos 14 cursos oferecidos por Instituições Públicas na região Centro-Oeste ( 77,8 \% do total regional ), nota -se que tais cursos apresentaram o conceito 2 como modal, atribuídos a cinco cursos.
Quatro cursos receberam conceito 3, e apenas um curso ficou sem conceito ( SC ).
Aos cursos restantes, foram atribuídos os conceitos 1 e 5 ( um curso \textbf{cada} ) e conceito 4 ( três cursos ).
O Gráfico 5.2 apresenta a distribuição do Conceito Enade segundo a Categoria Administrativa da IES.
Os cursos em IES Públicas ( linha azul ) apresentam uma poligonal distribuição bimodal, uma vez que as maiores frequências se concentraram em dois conceitos, os conceitos 2 e 3, atribuídos, \textbf{cada} qual, a três cursos ( 42,9 \% \textbf{cada} ).
Além disso, um curso alcançou o conceito máximo, 5.
Nenhum curso ficou sem conceito ( SC ) ou com conceito 1 ou 4 nessa combinação de categoria e região.
Na região Nordeste, a rede privada participou com somente um dos 26 cursos participantes, o equivalente a 3,8 \% do total da região.
Este curso recebeu conceito 2.
As Instituições Públicas da região Nordeste concentraram 25 cursos ( 96,2 \% do total da região ).
Desses, dez obtiveram conceito 3, o valor modal, seguidos por outros seis cursos, que receberam conceito 2.
Os conceitos 1 e 4 foram atribuídos, respectivamente, a três e a cinco cursos.
Apenas um curso ficou sem conceito ( SC ) nessa região, enquanto nenhum recebeu conceito 5.
Na região Sudeste, a proporção de cursos da rede privada foi de 22,7 \%, correspondendo a cinco dos 22 cursos participantes.
Nessa categoria e região, o conceito modal foi 3, atribuído a dois cursos.
Os três cursos restantes foram alocados aos conceitos 2, 4 e 5 ( um curso, \textbf{cada} ).
Nenhum curso ficou sem conceito ( SC ) ou com conceito 1.
Entre os 17 cursos oferecidos por Instituições Públicas na região Sudeste, o conceito modal foi 3, atribuído a 11 cursos.
Os demais cursos receberam os conceitos 2 ( dois cursos ) e 4 ( quatro cursos ).
Nenhum curso ficou sem conceito ( SC ) ou recebeu conceito 1 ou 5.
As Instituições Públicas concentraram todos os 16 cursos participantes da região Sul, 100,0 \% do total regional, a maior das proporções desta categoria nas regiões.
Desses, oito obtiveram conceito 3, o conceito modal.
Nesta combinação de Categoria Administrativa e Grande Região, os conceitos 2 e 4 foram atribuídos a dois e a cinco cursos, respectivamente.
Um curso recebeu conceito 5, e nenhum recebeu conceito 1 ou ficou sem conceito ( SC ).
As Instituições Públicas na região Sul não tiveram participação.
Na região Centro-Oeste, quatro dos 18 cursos participantes eram de Instituições Privadas ( 22,2 \% em termos regionais ), todos eles receberam conceito 2.
Quanto aos 14 cursos oferecidos por Instituições Públicas na região Centro-Oeste ( 77,8 \% do total regional ), nota -se que tais cursos apresentaram o conceito 2 como modal, atribuídos a cinco cursos.
Quatro cursos receberam conceito 3, e apenas um curso ficou sem conceito ( SC ).
Aos cursos restantes, foram atribuídos os conceitos 1 e 5 ( um curso \textbf{cada} ) e conceito 4 ( três cursos ).
O Gráfico 5.2 apresenta a distribuição do Conceito Enade segundo a Categoria Administrativa da IES.
Os cursos em IES Públicas ( linha azul ) apresentam uma poligonal

% matched lemmas: cada
\end{textsample}
